\metatitle{Operations and minors}
\metadescription{Loops, coloops, restrictions, deletion, contraction, and
direct sums of colored matroids.}

\section[operations]{Operations and minors}

The usual matroid operations extend to colored matroids. We define loops and
coloops one color at a time, while deletion and contraction involve both a
position and a color.

\subsection[loopsAndColoops]{Loops and coloops}

Let $\mathcal F$ be a colored matroid on $E$. For each color $c$, the sets
$B(u,c)$, where $u\in\mathcal F$, are the bases of a matroid on $E$. We write
$M(\mathcal F,c)$ for this color projection.

\begin{definition}[Loops and coloops]
A position $e\in E$ is a \defin{loop of color $c$} if no word in
$\mathcal F$ assigns color $c$ to $e$. It is a \defin{coloop of color $c$}
if every word in $\mathcal F$ assigns color $c$ to $e$. Equivalently, $e$ is
a loop or coloop, respectively, of $M(\mathcal F,c)$.
\end{definition}

\begin{example}
Consider the seven-word colored matroid
\[
 \mathcal F=\left\{\mathtt{1123},\mathtt{1213},\mathtt{1321},
 \mathtt{2113},\mathtt{3121},\mathtt{2311},\mathtt{3211}\right\}.
\]
Color $2$ never occurs at position $4$, so position $4$ is a loop of color
$2$. Similarly, position $3$ is a loop of color $3$.

If $\mathcal G=\{\mathtt{1231},\mathtt{1321}\}$, then positions $1$ and $4$
are coloops of color $1$. At either position, colors $2$ and $3$ are loops.
\end{example}

\subsection[restrictionsAndPinning]{Restrictions and pinning}

\begin{definition}[Restrictions and pinning]
Choose a set $L_e$ of allowed colors at each position $e$. The corresponding
\defin{restriction} of $\mathcal F$ is
\[
 \mathcal F_L
 \coloneqq
 \{u\in\mathcal F:u(e)\in L_e\text{ for every }e\in E\}.
\]
To \defin{pin} $(e,c)$ means to retain only the words with $u(e)=c$.
\end{definition}

\begin{proposition}[Restrictions]
Every nonempty restriction of a colored matroid is a colored matroid. In
particular, every nonempty pinned family satisfies symmetric cycle exchange.
\end{proposition}

\begin{proof}
Take two words in the restriction. Every entry used by a cyclic exchange is
an entry of one of the two words at the same position, so the two resulting
words still satisfy all allowed-color conditions. If a position is pinned,
the two words agree there, and the position does not occur in an exchange.
\end{proof}

\subsection[deletionAndContraction]{Deletion and contraction}

Removing a position without fixing its color would give words with different
quota vectors. We therefore define minors for a position-color pair $(e,c)$.

\begin{definition}[Deletion and contraction]
The \defin{deletion of $(e,c)$} is
\[
 \mathcal F\setminus(e,c)
 \coloneqq
 \{u\in\mathcal F:u(e)\ne c\}.
\]
The position $e$ remains, but color $c$ is no longer allowed there.

The \defin{contraction of $(e,c)$} is
\[
 \mathcal F/(e,c)
 \coloneqq
 \{u|_{E\setminus\{e\}}:u\in\mathcal F,\ u(e)=c\}.
\]
Contraction removes $e$ and decreases the quota of color $c$ by one.
\end{definition}

Both operations may give the empty family. If $e$ is a loop of color $c$,
then $\mathcal F/(e,c)$ is empty and
$\mathcal F\setminus(e,c)=\mathcal F$. If $e$ is a coloop of color $c$,
then $\mathcal F\setminus(e,c)$ is empty.

\begin{proposition}[Closure under deletion and contraction]
Every nonempty deletion or contraction of a colored matroid is a colored
matroid.
\end{proposition}

\begin{proof}
Deletion is the restriction that forbids color $c$ at $e$. Contraction first
pins $e$ to color $c$ and then removes the fixed position. Restrictions and
pinning preserve symmetric cycle exchange.
\end{proof}

\begin{example}[Contractions at one position]
For the seven-word family above, contraction at the first position gives
\[
\begin{array}{c|c}
c&\mathcal F/(1,c)\\ \hline
1&\{\mathtt{123},\mathtt{213},\mathtt{321}\}\\
2&\{\mathtt{113},\mathtt{311}\}\\
3&\{\mathtt{121},\mathtt{211}\}
\end{array}
\]
The three rows partition $\mathcal F$ according to the color at position $1$.
For example,
\[
 \mathcal F\setminus(1,1)
 =\{\mathtt{2113},\mathtt{3121},\mathtt{2311},\mathtt{3211}\}.
\]
\end{example}

\begin{proposition}[Ordinary matroid minors]
Suppose $\mathcal F$ has two colors, with color-$1$ classes the bases of a
matroid $M$. If $e$ is not a loop of $M$, then
$\mathcal F/(e,1)$ records the bases of $M/e$. If $e$ is not a coloop of
$M$, then $\mathcal F/(e,0)$ records the bases of $M\setminus e$.
\end{proposition}

\begin{proof}
The color-$1$ classes in $\mathcal F/(e,1)$ are the sets
$B\setminus\{e\}$ for bases $B$ containing $e$. These are the bases of
$M/e$. In $\mathcal F/(e,0)$, the color-$1$ classes are the bases of $M$
that avoid $e$, which are the bases of $M\setminus e$.
\end{proof}

\subsection[directSums]{Direct sums}

Let $\mathcal F_1$ and $\mathcal F_2$ use the same colors on disjoint
position sets $E_1$ and $E_2$.

\begin{definition}[Direct sum]
The \defin{direct sum} is
\[
 \mathcal F_1\oplus\mathcal F_2
 \coloneqq
 \{u_1\sqcup u_2:u_1\in\mathcal F_1,\ u_2\in\mathcal F_2\}.
\]
Its quota for each color is the sum of the two corresponding quotas.
\end{definition}

\begin{proposition}[Direct sums]
The direct sum $\mathcal F_1\oplus\mathcal F_2$ is a colored matroid if and
only if both factors are colored matroids. For every color $c$,
\[
 M(\mathcal F_1\oplus\mathcal F_2,c)
 =M(\mathcal F_1,c)\oplus M(\mathcal F_2,c).
\]
On the right, the bases are the unions of one basis from each projection.
\end{proposition}

\begin{proof}
Suppose first that both factors are colored matroids. A prescribed
disagreement belongs to one factor. Apply symmetric cycle exchange there and
leave the other factor unchanged.

Conversely, fix a word in one factor and compare two words that differ only
in the other factor. Every exchange between them stays in the positions where
they differ, and therefore gives the required exchange in that factor. The
statement about color projections follows from the definition.
\end{proof}

\begin{example}
Take
\[
 \mathcal F_1=\{\mathtt{123},\mathtt{231}\},
 \qquad
 \mathcal F_2=\{\mathtt{12},\mathtt{21}\}.
\]
Their direct sum is
\[
 \mathcal F_1\oplus\mathcal F_2
 =\{\mathtt{12312},\mathtt{12321},
     \mathtt{23112},\mathtt{23121}\},
\]
the five-position colored matroid from the introduction.
\end{example}
